% Originally: new section, 2026-08-09. Not a move --- this material had no home in the document
% before. Exercise 3.9 asked the reader to prove that all norms on R^n are equivalent, and the
% solution in 99-solutions.tex proved it, but the statement itself was nowhere in the body, so
% nothing could \cref it and the reader met the result only as homework.
% Supplement: lec_notes.pdf, pp. 31--33 (Definition 9.105, Theorem 9.107, Corollary 9.109)
% Generator: Claude Opus 5 (High)

\section{Equivalence of norms}
\label{sec:equivalence_of_norms}

The previous sections put several norms on the same vector space: the Euclidean norm, the
$p$-norms, the supremum norm. A natural worry follows. Convergence, openness, continuity and
compactness were all defined from a metric, and on a normed space that metric came from the norm.
So does changing the norm change which sequences converge and which sets are open?

On $\mathbb{R}^n$ the answer is no, and the reason is worth isolating, because it is one of the
few places in this course where finite-dimensionality is doing essential work.

% Supplement: lec_notes.pdf, p. 31 (Definition 9.105)
\begin{definition}[Equivalent norms]
\label{def:equivalent_norms}
Let $V$ be a vector space over $\mathbb{R}$ and let $\|\cdot\|_1$ and $\|\cdot\|_2$ be two norms
on $V$. They are called \newterm{equivalent} \germanfor{äquivalent}, or \emph{comparable}, if
there are constants $A, B > 0$ with
\[\|v\|_1 \leq A\|v\|_2 \qquad\text{and}\qquad \|v\|_2 \leq B\|v\|_1 \qquad\text{for all } v \in V.\]
\end{definition}

\begin{remark}[Why two inequalities and not one]
A single inequality $\|v\|_1 \leq A\|v\|_2$ says only that $\|\cdot\|_2$-convergence implies
$\|\cdot\|_1$-convergence: it makes the first norm the coarser of the two. Equivalence is the
two-sided statement, and it is exactly what is needed for the two norms to have the same null
sequences, and hence the same convergent sequences, open sets and compact sets. That equivalence
of norms really is an equivalence relation is part 5 of \cref{ex:3.9}.
\end{remark}

% Supplement: lec_notes.pdf, p. 32 (Example 9.106)
\begin{example}[The $1$-norm against the maximum norm]
\label{ex:one_norm_vs_max_norm}
On $\mathbb{R}^n$,
\[\|v\|_\infty \leq \|v\|_1 \qquad\text{and}\qquad \|v\|_1 \leq n\|v\|_\infty,\]
so the two are equivalent with $A = 1$ and $B = n$. The left inequality holds because the largest
$|v_i|$ is one of the terms of the sum $\sum_i|v_i|$; the right because each of the $n$ terms of
that sum is at most $\max_i|v_i|$.

The constant $n$ on the right is sharp, as $v = (1,1,\dots,1)$ shows, and it is the first sign of
what goes wrong later: the constants in \cref{thm:all_norms_equivalent} are allowed to depend on
the dimension, and there is no dimension to depend on in an infinite-dimensional space.
\end{example}

% Supplement: lec_notes.pdf, p. 32 (Theorem 9.107)
\begin{theorem}[All norms on $\mathbb{R}^n$ are equivalent]
\label{thm:all_norms_equivalent}
Any two norms on $\mathbb{R}^n$ are equivalent in the sense of \cref{def:equivalent_norms}.
\end{theorem}

\begin{proof}
This is \cref{ex:3.9}, whose five parts are exactly the five steps of the argument, and it is
proved in full in the solutions to that exercise. In outline, with $|\cdot|$ the Euclidean norm
and $f$ an arbitrary norm: expanding $x$ in the standard basis and applying the triangle
inequality, homogeneity and Cauchy--Schwarz gives $f(x) \leq C_1|x|$; that bound makes $f$
Lipschitz, hence continuous, for the Euclidean metric; Weierstrass then gives $f$ a positive
minimum $c_2$ on the Euclidean unit sphere, which is compact by Heine--Borel
(\cref{thm:heine_borel}); and homogeneity spreads that bound from the sphere to all of
$\mathbb{R}^n$, giving $f(x) \geq c_2|x|$.
\end{proof}

\begin{remark}[Where each hypothesis is used]
\label{rem:where_finite_dimension_enters}
Both halves of the proof use finite-dimensionality, and they use it differently. The first
inequality expands $x$ in a \emph{finite} basis, so the sum of $n$ terms it produces is a sum and
not a series. The second applies Weierstrass on the unit sphere, which needs that sphere to be
\emph{compact}, and by Heine--Borel it is, in $\mathbb{R}^n$ and nowhere else.
\Cref{rem:norm_equivalence_needs_finite_dimension} exhibits the failure in
$C^0([0,1],\mathbb{R})$, and \cref{ex:polynomials_infinite_dimensional} gives the algebraically
simplest infinite-dimensional space to have in mind.
\end{remark}

% Supplement: lec_notes.pdf, p. 33 (9.108)
\begin{corollary}[Any finite-dimensional real vector space]
\label{cor:all_norms_equivalent_finite_dim}
Let $V$ be a real vector space of finite dimension $n$. Then any two norms on $V$ are equivalent.
\end{corollary}

\begin{proof}
Fix a basis $\mathcal{B} = \{e_1,\dots,e_n\}$ of $V$. The coordinate map
$\iota_{\mathcal{B}} : \mathbb{R}^n \to V$, $(x_1,\dots,x_n) \mapsto x_1e_1+\dots+x_ne_n$, is a
linear isomorphism, so a norm $\|\cdot\|$ on $V$ pulls back to the norm
$x \mapsto \|\iota_{\mathcal{B}}(x)\|$ on $\mathbb{R}^n$, and the two-sided inequalities of
\cref{def:equivalent_norms} are preserved in both directions under this identification. Applying
\cref{thm:all_norms_equivalent} on $\mathbb{R}^n$ and transporting back gives the claim.
\end{proof}

% Supplement: lec_notes.pdf, p. 33 (Corollary 9.109)
\begin{corollary}[All norms give the same topology]
\label{cor:norms_same_topology}
Topological properties of a subset $E \subseteq \mathbb{R}^n$ (openness, closedness, convergence
of sequences, compactness, connectedness) do not depend on which norm was used to define the
metric.
\end{corollary}

\begin{proof}
Let $d_1, d_2$ be the metrics induced by two norms. By \cref{thm:all_norms_equivalent} there is a
constant $C > 0$ with $C^{-1}d_1(x,y) \leq d_2(x,y) \leq Cd_1(x,y)$ for all $x,y$, since
$d_i(x,y) = \|x-y\|_i$ and the norms are equivalent. Hence every $d_1$-ball contains a
$d_2$-ball about the same centre and conversely, which by \cref{def:open_closed} makes a set open
for $d_1$ exactly when it is open for $d_2$. The two metrics therefore generate the same
topology, and every property defined purely in terms of the open sets agrees.
\end{proof}

\begin{ainote}[What this licenses, and what it does not]
\label{ainote:norm_equivalence_scope}
\Cref{cor:norms_same_topology} is the reason this document is free to say \qt{open}, \qt{compact}
or \qt{convergent} in $\mathbb{R}^n$ without ever naming a norm, and it is worth knowing that the
licence has been earned rather than assumed.

Three things it does \emph{not} give you.
\begin{itemize}
  \item It is about \emph{topological} properties only. Metric quantities are not preserved:
    the unit ball genuinely changes shape with the norm (see the three drawn in
    \cref{sec:unit_balls}), distances change, and so do Lipschitz constants. A $1$-Lipschitz map
    for one norm need not be $1$-Lipschitz for another, only Lipschitz with some other constant.
  \item It says nothing in infinite dimensions, where the conclusion is false
    (\cref{rem:where_finite_dimension_enters}).
  \item It does not make the choice of norm irrelevant in practice. Estimates are still written
    in whichever norm makes the constants come out cleanly, which is why the supremum norm is the
    one used on $C(K,\mathbb{R}^m)$ in \cref{sec:ascoli_arzela} even though that space is
    infinite-dimensional and the theorem above does not reach it.
\end{itemize}
\end{ainote}
